% ID: SEL-D0009-BD-E08
% ESTADO: cotejada. Segunda pasada visual completa y borrador compilado revisados.
% Fuente: Selectividad/matematicas_aplicadas/ORD_RESER_Modelo-3-Enunciados (Mat-Ap-2021).pdf
% SHA-256: 1dc0455361ce91a1157a7b3e4a06db95bc988e8dd0c9756e9aeb74b8cbb73a81
% Página: 2; bloque D; ejercicio original 8.

Se quiere estimar la proporción de imprentas de una región que incluyen
el uso de celulosa reciclada en los libros que imprimen. Para ello, se ha
tomado una muestra aleatoria de 50 imprentas de esa región y en ella hay
12 que usan dicho material.
\begin{enumerate}[label=\alph*)]
  \item \textbf{(1.5 puntos)} Obtenga un intervalo de confianza al 95\%,
  para estimar la proporción real de imprentas que usan celulosa reciclada.
  \item \textbf{(1 punto)} Determine el tamaño mínimo de la muestra de
  imprentas de esa región que se deben seleccionar para que, manteniendo
  el mismo nivel de confianza y proporción muestral anteriores, la amplitud
  del intervalo sea como máximo de 0.2.
\end{enumerate}
